\documentclass[11pt,a4paper]{article}

% =====================================================
% preamble.tex — Configuración común de estilo
% Para enunciados de ejercitación en Tekla Structures
% =====================================================

% --- Codificación e idioma ---
\usepackage[utf8]{inputenc}
\usepackage[T1]{fontenc}
\usepackage[spanish,es-tabla,es-noshorthands]{babel}
\usepackage{lmodern}
\usepackage{microtype}

% --- Geometría de página ---
\usepackage[a4paper,margin=2.5cm,headheight=14pt]{geometry}

% --- Matemática y unidades ---
\usepackage{amsmath,amssymb}
\usepackage{siunitx}
\sisetup{
  output-decimal-marker={,},
  per-mode=symbol,
  group-separator={.},
  group-minimum-digits=4
}

% --- Tablas ---
\usepackage{booktabs}
\usepackage{array}
\usepackage{multirow}
\usepackage{tabularx}
\usepackage{longtable}

% --- Figuras ---
\usepackage{graphicx}
\usepackage{caption}
\usepackage{subcaption}
\usepackage{float}

% --- Gráficos vectoriales ---
\usepackage{tikz}
\usetikzlibrary{arrows.meta,positioning,calc,patterns,shapes.geometric}

% --- Listas ---
\usepackage{enumitem}
\setlist{itemsep=2pt,parsep=2pt}

% --- Color y cajas ---
\usepackage{xcolor}
\definecolor{teklablue}{RGB}{0,87,156}
\definecolor{boxgray}{RGB}{245,245,245}
\definecolor{accentorange}{RGB}{230,120,30}

\usepackage[most]{tcolorbox}

% Caja para comandos / atajos
\newtcolorbox{comando}{
  colback=boxgray,
  colframe=teklablue,
  boxrule=0.5pt,
  arc=2pt,
  left=6pt, right=6pt, top=4pt, bottom=4pt,
  fonttitle=\bfseries,
  title=Comando Tekla
}

% Caja para tips / consejos
\newtcolorbox{tip}{
  colback=yellow!8,
  colframe=accentorange,
  boxrule=0.5pt,
  arc=2pt,
  fonttitle=\bfseries,
  title=Tip
}

% Caja para entregable
\newtcolorbox{entregable}{
  colback=green!5,
  colframe=green!50!black,
  boxrule=0.5pt,
  arc=2pt,
  fonttitle=\bfseries,
  title=Entregable
}

% Caja para advertencia
\newtcolorbox{advertencia}{
  colback=red!5,
  colframe=red!70!black,
  boxrule=0.5pt,
  arc=2pt,
  fonttitle=\bfseries,
  title=Atención
}

% --- Código y comandos en línea ---
\usepackage{listings}
\lstset{
  basicstyle=\ttfamily\small,
  breaklines=true,
  frame=single,
  backgroundcolor=\color{boxgray},
  rulecolor=\color{teklablue}
}

% Comando para resaltar comandos/menús de Tekla en línea
% Uso: \tk{File > Open}  o  \tk{Component Catalog}
\newcommand{\tk}[1]{\textsf{\textbf{\textcolor{teklablue}{#1}}}}

% Comando para nombres de archivo / rutas
\newcommand{\arch}[1]{\texttt{#1}}

% Comando para teclas / shortcuts
% Uso: \key{Ctrl+S}
\newcommand{\key}[1]{\fbox{\footnotesize\texttt{#1}}}

% --- Encabezado y pie de página ---
\usepackage{fancyhdr}
\pagestyle{fancy}
\fancyhf{}
\fancyhead[L]{\small\textit{\tituloEjercicio}}
\fancyhead[R]{\small\thepage}
\fancyfoot[C]{\small Tekla Structures 2022 — Nivel avanzado}
\renewcommand{\headrulewidth}{0.4pt}
\renewcommand{\footrulewidth}{0.4pt}

% --- Hipervínculos y referencias ---
\usepackage{hyperref}
\hypersetup{
  colorlinks=true,
  linkcolor=teklablue,
  urlcolor=teklablue,
  citecolor=teklablue,
  pdfborder={0 0 0}
}
\usepackage[noabbrev,spanish]{cleveref}

% --- Formato de secciones ---
\usepackage{titlesec}
\titleformat{\section}
  {\Large\bfseries\color{teklablue}}
  {\thesection.}{0.5em}{}
\titleformat{\subsection}
  {\large\bfseries\color{teklablue!80!black}}
  {\thesubsection.}{0.5em}{}

% =====================================================
% commands.tex — Variables propias de cada enunciado
% Editar estos campos al inicio para cada ejercicio
% =====================================================

% --- Identificación del ejercicio ---
\newcommand{\tituloEjercicio}{Ejercicio: [TÍTULO]}
\newcommand{\subtituloEjercicio}{[Subtítulo descriptivo]}
\newcommand{\codigoEjercicio}{TEKLA-XXX}
\newcommand{\nivelEjercicio}{Avanzado}
\newcommand{\duracionEjercicio}{[X] horas}

% --- Datos del curso ---
\newcommand{\nombreCurso}{[Nombre del curso]}
\newcommand{\autorEnunciado}{[Docente]}
\newcommand{\fechaEjercicio}{\today}

% --- Versión del software ---
\newcommand{\versionTekla}{Tekla Structures 2022}
\newcommand{\entornoTekla}{[Environment / Configuration]}

% --- Atajos para terminología recurrente ---
\newcommand{\CC}{\tk{Component Catalog}}
\newcommand{\AC}{\tk{Applications \& Components}}
\newcommand{\OB}{\tk{Object Browser}}
\newcommand{\OR}{\tk{Object Representation}}
\newcommand{\OG}{\tk{Object Group}}
\newcommand{\SF}{\tk{Selection Filter}}
\newcommand{\VF}{\tk{View Filter}}


\renewcommand{\tituloEjercicio}{Ejercicio de Estructura Metálica}
\renewcommand{\subtituloEjercicio}{[Subtítulo: tipo de estructura]}
\renewcommand{\codigoEjercicio}{TEKLA-EM-01}

\begin{document}

% =====================================================
% PORTADA
% =====================================================
\begin{titlepage}
\centering
\vspace*{2cm}
{\Huge\bfseries\color{teklablue} \tituloEjercicio \par}
\vspace{0.5cm}
{\Large \subtituloEjercicio \par}
\vspace{2cm}
\rule{0.6\textwidth}{0.4pt}\par
\vspace{1cm}
\begin{tabular}{rl}
\textbf{Código:} & \codigoEjercicio \\
\textbf{Curso:} & \nombreCurso \\
\textbf{Nivel:} & \nivelEjercicio \\
\textbf{Duración estimada:} & \duracionEjercicio \\
\textbf{Software:} & \versionTekla \\
\textbf{Docente:} & \autorEnunciado \\
\textbf{Fecha:} & \fechaEjercicio \\
\end{tabular}
\vfill
\end{titlepage}

\tableofcontents
\newpage

% =====================================================
\section{Objetivos de aprendizaje}
% =====================================================
\begin{itemize}
    \item [Objetivo 1: ej. modelar elementos con \tk{Steel Beam}, \tk{Steel Column}, \tk{Steel Plate}]
    \item [Objetivo 2: ej. aplicar conexiones del \CC]
    \item [Objetivo 3: ej. crear un \tk{Custom Component} para una unión no estándar]
    \item [Objetivo 4: ej. trabajar con \tk{Assemblies} y \tk{Sub-Assemblies}]
    \item [Objetivo 5: ej. producir planos de taller y lista de bulones]
\end{itemize}

% =====================================================
\section{Descripción del proyecto}
% =====================================================
[Descripción breve del proyecto y contexto profesional.]

% =====================================================
\section{Datos de partida}
% =====================================================
\subsection{Geometría general}
\begin{itemize}
    \item [Luces, alturas, cantidad de pórticos]
    \item [Cargas o solicitaciones de referencia (si aplica)]
\end{itemize}

\subsection{Perfiles y materiales}
\begin{table}[H]
\centering
\begin{tabular}{lll}
\toprule
\textbf{Elemento} & \textbf{Perfil} & \textbf{Acero} \\
\midrule
[Columnas]        & [HEB / W]   & [F-24 / A572] \\
[Vigas principales]& [IPE / W]   & [F-24] \\
[Vigas secundarias]& [UPN / C]   & [F-24] \\
[Arriostramientos]& [L / 2L]    & [F-24] \\
[Chapas]          & [PL XX]     & [F-24] \\
\bottomrule
\end{tabular}
\caption{Perfiles y aceros del proyecto.}
\end{table}

\subsection{Bulones}
\begin{itemize}
    \item Calidad: [A325 / 8.8 / etc.]
    \item Diámetros principales: [M16, M20, M24]
\end{itemize}

\subsection{Normativa de referencia}
\begin{itemize}
    \item [CIRSOC 301 / 304 / AISC 360 / EN 1993]
    \item [Estándar de planos de taller]
\end{itemize}

% =====================================================
\section{Esquema de referencia}
% =====================================================
\begin{figure}[H]
\centering
\fbox{\parbox[c][6cm][c]{0.8\textwidth}{\centering\textit{[Esquema de referencia: planta, cortes y elevaciones]}}}
\caption{Esquema general del proyecto.}
\end{figure}

% =====================================================
\section{Consigna paso a paso}
% =====================================================

\subsection{Configuración inicial}
\begin{enumerate}
    \item Crear modelo nuevo con el entorno \entornoTekla.
    \item Configurar \tk{Project Properties}.
    \item Definir \tk{Grids} y \tk{Levels}.
    \item Configurar perfiles y materiales en \tk{Catalogs}.
\end{enumerate}

\subsection{Modelado de la estructura primaria}
\begin{enumerate}
    \item Modelar columnas con \tk{Steel Column}.
    \item Modelar vigas principales con \tk{Steel Beam}.
    \item Modelar arriostramientos.
    \item Verificar continuidad y conexiones lógicas.
\end{enumerate}

\subsection{Conexiones estándar}
\begin{enumerate}
    \item Aplicar conexiones desde el \CC.
    \item Configurar parámetros de cada conexión: [tipo, espesores, bulones].
    \item Verificar con \tk{Inquire Component} y revisar warnings.
\end{enumerate}

\begin{comando}
[Detalle de un componente clave: número y configuración recomendada]
\end{comando}

\subsection{Custom Component}
\begin{enumerate}
    \item Crear un \tk{Custom Component} tipo \textit{Connection / Detail / Part} para [unión específica].
    \item Definir variables paramétricas en \tk{Custom Component Editor}.
    \item Configurar \tk{Distance Variables} y \tk{Reference Models}.
    \item Probar el componente en al menos 3 casos distintos para validar paramétricamente.
\end{enumerate}

\begin{tip}
[Tip sobre cómo estructurar las variables del custom component]
\end{tip}

\subsection{Numeración y assemblies}
\begin{enumerate}
    \item Configurar \tk{Numbering Settings}.
    \item Ejecutar \tk{Numbering > Number Modified Objects}.
    \item Revisar \tk{Assemblies} y \tk{Sub-Assemblies} resultantes.
\end{enumerate}

\subsection{Planos de taller y reportes}
\begin{enumerate}
    \item Generar \tk{General Arrangement Drawings}.
    \item Generar \tk{Assembly Drawings} de los conjuntos clave.
    \item Generar \tk{Single Part Drawings} de piezas representativas.
    \item Producir reportes: lista de materiales, lista de bulones, peso total.
\end{enumerate}

% =====================================================
\section{Entregables}
% =====================================================
\begin{entregable}
\begin{enumerate}
    \item Archivo de modelo \arch{.db1} comprimido.
    \item \tk{Custom Component} exportado como \arch{.uel}.
    \item Planos de conjunto (assembly drawings) en PDF.
    \item Lista de materiales y lista de bulones.
    \item Documento con capturas de hitos clave.
\end{enumerate}
\end{entregable}

% =====================================================
\section{Criterios de evaluación}
% =====================================================
\begin{table}[H]
\centering
\begin{tabularx}{\textwidth}{Xcc}
\toprule
\textbf{Criterio} & \textbf{Puntaje} & \textbf{Obtenido} \\
\midrule
Geometría y perfiles correctos & 15 & \\
Aplicación de conexiones estándar & 20 & \\
Calidad y robustez del custom component & 25 & \\
Coherencia de numeración y assemblies & 10 & \\
Calidad de planos de taller & 20 & \\
Completitud de entregables & 10 & \\
\midrule
\textbf{Total} & \textbf{100} & \\
\bottomrule
\end{tabularx}
\end{table}

% =====================================================
\section{Tips y errores comunes}
% =====================================================
\begin{tip}
[Tip 1]
\end{tip}

\begin{advertencia}
[Error frecuente]
\end{advertencia}

% =====================================================
\section{Referencias y documentación}
% =====================================================
\begin{itemize}
    \item Tekla User Assistance: \url{https://support.tekla.com/}
    \item [Otras referencias]
\end{itemize}

\end{document}
